% =============================================================================
%  导言区：中文译注本《给数学家的导出范畴》
%  原文：R. P. Thomas, "Derived categories for the working mathematician",
%        arXiv:math/0001045（v1 2000-01-07，v2 2001-10-17），13 pp.
%        会议版：Winter School on Mirror Symmetry, Vector Bundles and
%        Lagrangian Submanifolds（Harvard, 1999），
%        AMS/IP Stud. Adv. Math. 23, 349--361 (2001).
%  编译：xelatex main.tex（两遍）
% =============================================================================
\documentclass[11pt,a4paper]{article}

\usepackage[UTF8, fontset=fandol, heading=true]{ctex}
\usepackage{amsmath}
\usepackage{amssymb}
\usepackage{amsthm}
\usepackage{mathtools}
\usepackage{bm}
\usepackage[a4paper,left=2.5cm,right=2.5cm,top=2.6cm,bottom=2.6cm]{geometry}
\usepackage{graphicx}
\usepackage{xcolor}
\usepackage[most]{tcolorbox}
\usepackage{enumitem}
\usepackage{booktabs}
\usepackage{longtable}
\usepackage{array}
\usepackage{tikz}
\usepackage{tikz-cd}
\usetikzlibrary{shapes.geometric,arrows.meta,positioning,calc,decorations.pathmorphing}
\usepackage[font=small,labelfont=bf,labelsep=quad]{caption}
\usepackage{fancyhdr}
\usepackage[colorlinks=true,
            linkcolor=blue!45!black,
            citecolor=blue!45!black,
            urlcolor=blue!45!black]{hyperref}

\linespread{1.15}
\setlength{\emergencystretch}{3em}
\allowdisplaybreaks[3]
\renewcommand{\figurename}{图}
\renewcommand{\tablename}{表}
\renewcommand{\refname}{参考文献}
\renewcommand{\contentsname}{目录}
\renewcommand{\appendixname}{附录}
\renewcommand{\proofname}{证明}
\renewcommand{\abstractname}{摘要}

% 页眉页脚：双面排印时左页书名、右页节名
\pagestyle{fancy}
\fancyhf{}
\fancyhead[L]{\small\itshape 给数学家的导出范畴}
\fancyhead[R]{\small\itshape\leftmark}
\fancyfoot[C]{\small\thepage}
\renewcommand{\headrulewidth}{0.4pt}
\fancypagestyle{plain}{\fancyhf{}\fancyfoot[C]{\small\thepage}%
  \renewcommand{\headrulewidth}{0pt}}

% ---------------------------------------------------------------------------
% 原文字体约定：原稿把 \emph 重定义为无衬线粗体（amsproc + 作者宏），
% 中译本改作粗体，用于标记作者强调的关键术语。
% ---------------------------------------------------------------------------
\renewcommand{\emph}[1]{{\bfseries #1}}

% ---------------------------------------------------------------------------
% 原文数学宏（与原文 preamble 逐一对应，保证公式可逐符号对照）
%   \Z \C \R \OO \A \B \F \RR \Hom \udot \ldot \Ltimes \ve \res \To \into
% ---------------------------------------------------------------------------
\newcommand{\Z}{\mathbb{Z}}
\newcommand{\C}{\mathbb{C}}
\newcommand{\R}{\mathbf{R}}
\newcommand{\OO}{\mathcal{O}}
\newcommand{\A}{\mathcal{A}}
\newcommand{\B}{\mathcal{B}}
\newcommand{\F}{\mathcal{F}}
\newcommand{\RR}{\mathcal{R}}
\newcommand{\Hom}{\operatorname{Hom}}
\newcommand{\Ext}{\operatorname{Ext}}
\newcommand{\SheafExt}{\mathcal{E}xt}
\newcommand{\Tor}{\operatorname{Tor}}
\newcommand{\Spec}{\operatorname{Spec}}
\newcommand{\cyl}{\operatorname{cyl}}
\newcommand{\id}{\operatorname{id}}
\newcommand{\into}{\hookrightarrow}
\newcommand{\To}{\longrightarrow}
\newcommand{\ot}{\leftarrow}
\newcommand{\oT}{\longleftarrow}
\newcommand{\udot}{^{\scriptscriptstyle\bullet}}
\newcommand{\ldot}{_{\scriptscriptstyle\bullet}}
\newcommand{\Ltimes}{\mathbin{\buildrel{\mathbf L}\over{\otimes}}}
\newcommand{\ve}{^{\vee}}
\newcommand{\res}{\arrowvert}
\newcommand{\Aut}{\operatorname{Aut}}

% ---------------------------------------------------------------------------
% 定理类环境：与原文一致——共享一个计数器、按节编号。
% 原文仅用 theorem/definition/eg/remark 四种，中译本补上引理与命题以备引用。
% ---------------------------------------------------------------------------
\newtheoremstyle{cnth}%
  {6pt}% above
  {6pt}% below
  {\normalfont}% body
  {}% indent
  {\bfseries}% head font
  {.}% punctuation
  {0.5em}% space after head
  {}% head spec

\theoremstyle{cnth}
\newtheorem{thm}{定理}[section]
\newtheorem{lemma}[thm]{引理}
\newtheorem{prop}[thm]{命题}
\newtheorem{cor}[thm]{推论}
\newtheorem{definition}[thm]{定义}
\newtheorem{eg}[thm]{例}
\newtheorem{remark}[thm]{注}

% ---------------------------------------------------------------------------
% 三类译者方框：yizhu（橙色，编号"译注 N"）、conceptbox（紫色，概念卡片）、
% keytakeaways（蓝色，本节重点）。三个 \ifshow* 开关可一键关闭。
% ---------------------------------------------------------------------------
\newif\ifshowyizhu
\showyizhutrue
\newcounter{yizhu}

\newenvironment{yizhu}[1][]{%
    \refstepcounter{yizhu}%
    \par\medskip
    \ifshowyizhu
        \begin{tcolorbox}[enhanced, breakable,
            colback=orange!4!white, colframe=orange!55!black,
            boxrule=0.5pt, arc=2pt, left=6pt, right=6pt, top=4pt, bottom=4pt,
            title={\small\textbf{译注 \theyizhu\if\relax\detokenize{#1}\relax\else\ ——#1\fi}},
            fonttitle=\bfseries, coltitle=black, colbacktitle=orange!12!white,
            toptitle=1pt, bottomtitle=1pt]
    \fi
}{%
    \ifshowyizhu
        \end{tcolorbox}%
    \fi
    \par\smallskip
}

\newif\ifshowconcept
\showconcepttrue

\newenvironment{conceptbox}[1]{%
    \par\medskip
    \ifshowconcept
        \begin{tcolorbox}[enhanced, breakable,
            colback=violet!3!white, colframe=violet!55!black,
            boxrule=0.5pt, arc=2pt, left=6pt, right=6pt, top=4pt, bottom=4pt,
            title={\small\textbf{概念卡片：#1}\quad{\footnotesize（译者所加）}},
            fonttitle=\bfseries, coltitle=black, colbacktitle=violet!12!white,
            toptitle=1pt, bottomtitle=1pt]
    \fi
}{%
    \ifshowconcept
        \end{tcolorbox}%
    \fi
    \par\medskip
}

\newif\ifshowkey
\showkeytrue

\newenvironment{keytakeaways}{%
    \par\medskip
    \ifshowkey
        \begin{tcolorbox}[enhanced, breakable,
            colback=blue!3!white, colframe=blue!35!black,
            boxrule=0.5pt, arc=2pt, left=6pt, right=6pt, top=4pt, bottom=4pt,
            title={\small\textbf{本节重点}\quad{\footnotesize（译者所加）}},
            fonttitle=\bfseries, coltitle=black, colbacktitle=blue!10!white,
            toptitle=1pt, bottomtitle=1pt]
    \fi
    \begin{itemize}[leftmargin=1.4em,itemsep=0.15em,topsep=2pt]
}{%
    \end{itemize}%
    \ifshowkey
        \end{tcolorbox}%
    \fi
    \par\medskip
}

% ---------------------------------------------------------------------------
% 校勘框：标出原文疑误之处，正文用脚注体例 N，详情见附录 B。
% ---------------------------------------------------------------------------
\newif\ifshowkanwu
\showkanwutrue
\newcounter{kanwu}

\newenvironment{kanwu}[1][]{%
    \refstepcounter{kanwu}%
    \par\medskip
    \ifshowkanwu
        \begin{tcolorbox}[enhanced, breakable,
            colback=red!3!white, colframe=red!50!black,
            boxrule=0.5pt, arc=2pt, left=6pt, right=6pt, top=4pt, bottom=4pt,
            title={\small\textbf{校勘 \thekanwu\if\relax\detokenize{#1}\relax\else\ ——#1\fi}},
            fonttitle=\bfseries, coltitle=black, colbacktitle=red!10!white,
            toptitle=1pt, bottomtitle=1pt]
    \fi
}{%
    \ifshowkanwu
        \end{tcolorbox}%
    \fi
    \par\smallskip
}

% 译注/校勘交叉标记
\newcommand{\yizhuRef}[1]{{\color{orange!70!black}\textsuperscript{译注~#1}}}
\newcommand{\kanwuRef}[1]{{\color{red!70!black}\textsuperscript{校勘~#1}}}

% 术语首次出现：中文（English）
\newcommand{\term}[2]{#1（\textit{#2}）}

% 引用文献的短命令（与原文 \cite 对应）
\newcommand{\cGM}{\cite{GM}}
\newcommand{\cHRD}{\cite{HRD}}
\newcommand{\cK}{\cite{K}}
\newcommand{\cMa}{\cite{Ma}}
